\begin{lemma}
\label{lemma-artinian-finite-length}
A ring $R$ is Artinian if and only if it has finite length
as a module over itself. Any such ring $R$ is both Artinian and
Noetherian, any prime ideal of $R$ is a maximal ideal, and $R$ is equal
to the (finite) product of its localizations at its maximal ideals.
\end{lemma}

\begin{proof}
If the original module $R$ has finite length, the proof of
Lemma \ref{lemma-finite-length-finite} gives both chain conditions
for its submodules, which are exactly its ideals. The descending
condition says that $R$ is Artinian. The ascending condition says
that every ideal is finite by the arbitrary-module equivalence
proved in Lemma \ref{lemma-Noetherian-basic}, so $R$ is Noetherian.

Conversely, suppose $R$ is Artinian. Lemmas
\ref{lemma-artinian-finite-nr-max},
\ref{lemma-artinian-radical-nilpotent}, and
\ref{lemma-product-local} identify $R$ with the finite product
of its original localizations $R_{\mathfrak m_i}$ and prove
that every prime is maximal. The product comparison is the
canonical map $r\mapsto(r/1)_i$, with the inverse proved there.
Its component maps are surjective: in the preceding proof each
localization is identified with the quotient $e_iR$ of $R$.
Thus each $R_{\mathfrak m_i}$ is Artinian, since a descending
chain of its ideals pulls back to such a chain in $R$.
When $R=0$, the product is empty and the length is zero.
For nonzero $R$, the module decomposition and length formula
in that same proof show that it suffices to prove finite length
for each original local factor.

Consider one such local Artinian ring $C=R_{\mathfrak m_i}$,
with maximal ideal $\mathfrak n=\mathfrak m_i C$ and residue
field $k=C/\mathfrak n$. By the nilpotence lemma,
$\mathfrak n^N=0$ for some $N\geq1$.
Each layer $V_j=\mathfrak n^j/\mathfrak n^{j+1}$,
$0\leq j<N$, is a $k$-vector space via the original quotient
map, because $\mathfrak n V_j=0$.
This space must have finite dimension. If not, choose a basis
containing distinct elements $v_1,v_2,\ldots$.
For each $a\geq0$ let $W_a$ be the span of all basis elements
except $v_1,\ldots,v_a$. Then
$W_0\supsetneq W_1\supsetneq\cdots$ is a strict descending
chain of vector subspaces. Their inverse images under
$\mathfrak n^j\to V_j$ are ideals of $C$ containing
$\mathfrak n^{j+1}$, and strictness is preserved by the
surjectivity of that quotient map. This contradicts the
Artinian condition on $C$.
By Lemma \ref{lemma-dimension-is-length}, the length of each
layer is its finite $k$-dimension. Additivity along the original
filtration gives
$$
\operatorname{length}_C(C)
=\sum_{j=0}^{N-1}\dim_k(\mathfrak n^j/\mathfrak n^{j+1})<\infty.
$$
Summing these finite lengths over the original local factors
proves finite length of $R$ and completes all assertions.

There are quantitative consequences for the same ring.
Write $\ell_i=\operatorname{length}_{R_{\mathfrak m_i}}
(R_{\mathfrak m_i})$. Each is positive and finite, and
$\operatorname{length}_R(R)=\sum_i\ell_i$.
The bound proved in Lemma \ref{lemma-length-infinite} gives
$(\mathfrak m_iR_{\mathfrak m_i})^{\ell_i}=0$.
Under the canonical product map, the original Jacobson radical
$I$ maps onto $\prod_i\mathfrak m_iR_{\mathfrak m_i}$:
membership in each component maximal ideal is exactly membership
in the corresponding original maximal ideal, and the product
map is an isomorphism. Powers are componentwise, so for
$R\neq0$,
$$
I^{\max_i\ell_i}=0.
$$
For $R=0$ the radical is already zero and no maximum of an empty
set is used.
Finally, every finite $R$-module $M$ has finite length. If its
given generating list has $r$ members, the associated surjection
$R^r\to M$ and additivity of length give the precise bound
$$
\operatorname{length}_R(M)\leq
r\operatorname{length}_R(R).
$$
The reverse implication, from finite length to finite generation,
is Lemma \ref{lemma-finite-length-finite}. Thus over this
Artinian ring the two conditions on $M$ are equivalent, with
the displayed bound retaining the chosen number of generators.
\end{proof}